\section{Overlap arrays and probability cascades}
\label{arr:section}

Fresh Gaussian patterns are determined on finitely many replicas by their
overlap covariance matrices.  This lets us evaluate their contribution
from limiting array laws, without representing an arbitrary limiting Gibbs
measure.  We first derive the ultrametric structure forced by the
Ghirlanda--Guerra identities and realize its finite approximations by
probability cascades.  We then compute how Gaussian marks change when a
cascade is reweighted.  Finally, Lemma~\ref{arr:limit} expresses the
reciprocal normalizers of bounded reweightings as limits of polynomials;
each polynomial uses only finitely many replicas.  Thus disorder-averaged
array laws suffice even for these normalized expectations.

All array laws in this section include the disorder average.  In particular,
the conditional overlap identities below condition on a sampled array
block, not on a fixed realization of the disorder.

For a Gibbs probability, replicas are independent samples conditional on its
reference data and disorder; their conditional expectation is denoted by
$\langle\cdot\rangle$.  The expectation $\mathbb E$ also averages the disorder
and any additional randomness specified in the construction.  A symmetric
array $B=(B_{\ell m})_{\ell,m\geq1}$ is \emph{weakly exchangeable} if its law is
invariant under every finite permutation of the indices, and is a
\emph{Gram array} if every finite restriction is positive semidefinite.
These properties, and a fixed diagonal, pass to limits in distribution of
finite restrictions.  Bounded arrays admit such subsequential limits.

\subsection{Conditional identities, positivity, and duplication}

The overlap-mixture identities originate in the work of Ghirlanda and
Guerra~\cite{gg1998}.  Their occurrence in the present models will be
proved by the Gaussian perturbations in the upper and lower bounds.

\begin{definition}[Ghirlanda--Guerra property]
\label{arr:gg}
Let $B$ be a weakly exchangeable array with values in a compact space, and
write $\mu=\mathcal L(B_{12})$.  We say that $B$ has the
Ghirlanda--Guerra property if, for every $n\geq2$,
\begin{equation}
 \mathcal L\bigl(B_{1,n+1}\mid (B_{\ell m})_{\ell,m\leq n}\bigr)
 =\frac1n\mu+\frac1n\sum_{\ell=2}^n\delta_{B_{1\ell}}.
 \label{arr:conditional-gg}
\end{equation}
Equivalently, for every bounded measurable function $F$ of the first
$n\times n$ block and every bounded measurable $\psi$,
\begin{equation}
 \mathbb E[F\psi(B_{1,n+1})]
 =\frac1n\mathbb E[F]\mathbb E[\psi(B_{12})]
  +\frac1n\sum_{\ell=2}^n\mathbb E[F\psi(B_{1\ell})].
 \label{arr:gg-tested}
\end{equation}
The same definition applies to a joint array, with $B_{\ell m}$ a tuple of
entries.  By exchangeability the distinguished index $1$ can be replaced
by any of the first $n$ indices.
\end{definition}

It suffices to establish \eqref{arr:gg-tested} for continuous $F,\psi$:
the equality then identifies the relevant finite measures, and hence holds
for bounded measurable tests.  This observation also proves that the
property passes to weak limits on a compact entry space.  A measurable
entrywise image of a joint array retains the property, with conditioning
on its own first block, by the conditional-expectation tower rule.

The direct array argument below uses duplication, the mechanism in
Panchenko's ultrametricity proof~\cite[Theorem~4 and Section~3]{panchenko2013}.

\begin{theorem}[Positivity and ultrametricity]
\label{arr:ultrametric}
Suppose $B$ is a bounded, symmetric, weakly exchangeable Gram array with
constant diagonal $D$ and the Ghirlanda--Guerra property.  Then, almost
surely, for distinct indices,
\begin{equation}
 B_{\ell m}\geq0,
 \qquad
 B_{12}\geq\min\{B_{13},B_{23}\}.
 \label{arr:ultrametric-inequality}
\end{equation}
\end{theorem}

\begin{proof}
First suppose $b=\mathbb P(B_{12}<-\varepsilon)>0$ for some
$\varepsilon>0$.  On the event that all distinct pairs among the first
$n$ indices have overlap below $-\varepsilon$, each of the $n$ links to
index $n+1$ fails this condition with conditional probability
$(1-b)/n$, by \eqref{arr:conditional-gg}.  A union bound therefore gives
conditional probability at least $b$ that all the new links satisfy it.
Starting with $n=2$, such blocks have positive probability for every size.
Their quadratic form on the all-ones vector is at most
$nD-n(n-1)\varepsilon$, contradicting positive semidefiniteness for large
$n$.  Taking a countable sequence of $\varepsilon$ proves positivity.

We next prove a support duplication assertion.  If
$(a_{\ell m})_{\ell,m\leq n}$ belongs to the support of a finite block,
$n\geq2$, there is a supported extension to $n+1$ such that
\begin{equation}
 a_{i,n+1}=a_{in}\quad(i<n),\qquad
 a_{n,n+1}\leq\max_{i<n}a_{in}.
 \label{arr:duplication}
\end{equation}
Fix $b_*>\max_{i<n}a_{in}$ and a positive-probability neighborhood event
$A$ for the first $n$ indices, chosen so that all specified links to $n$
are below $b_*$.  Let $A_0$ impose just the restrictions among the first
$n-1$ indices.  An index $j\geq n$ \emph{matches} if its links to those
$n-1$ indices belong to the neighborhoods prescribed for $n$.  Thus
$A=A_0\cap\{n\text{ matches}\}$.

Suppose that, on $A$, a fresh index cannot both match and have its link to
$n$ below $b_*$.  Let $W_B$ be the limiting empirical frequency of
matching indices, and $W$ the frequency of indices $j>n$ with
$B_{nj}\geq b_*$.  These frequencies exist in $L^2$.  Indeed, for an
exchangeable sequence of indicators the second moments of empirical
averages and their mixed second moments have the same limit, so the
averages are $L^2$-Cauchy.  The limits are unchanged by finite
permutations of the relevant unobserved indices.  The supposition and
exchangeability imply $W_B\leq W$ on $A$.

Put $b=\mathbb P(B_{12}\geq b_*)$.  We have $b<1$, since $A$ has positive
probability and specifies links below $b_*$.  Conditional on $A$, the
first high link to $n$ has probability $b/n$; after $j$ consecutive high
links the next has probability $(b+j)/(n+j)$.  Applying
\eqref{arr:conditional-gg} successively, and then expressing empirical
frequency moments by products at distinct fresh indices, gives
\begin{equation}
 \mathbb E[\mathbf1_A W^d]
 =\mathbb P(A)\prod_{j=0}^{d-1}\frac{b+j}{n+j},
 \qquad d\geq1.
 \label{arr:duplication-moments}
\end{equation}
The binomial series therefore yields
\begin{equation}
 \mathbb E[\mathbf1_A(1+zW)^{-n}]
 =\mathbb P(A)(1+z)^{-b},\qquad z>0.
 \label{arr:duplication-transform}
\end{equation}
Initially the series proves this for $0<z<1$.  Both sides are real
analytic on $(0,\infty)$, the left side by boundedness of $W$ and local
power-series domination, so the identity holds throughout that interval.
It also includes $b=0$.

To compare the matching frequency with the high-link frequency, write
$I_j=\mathbf1_{\{j\text{ matches}\}}$.  The event $A_0$ uses only the
first $n-1$ indices, and $W_B$ is unchanged by every finite permutation
of the remaining indices, including index $n$.  Hence, for every bounded
continuous $\varphi$, exchangeability gives
\[
 \mathbb E[\mathbf1_{A_0}I_n\varphi(W_B)]
 =\mathbb E\left[\mathbf1_{A_0}
   \frac1M\sum_{j=n}^{n+M-1}I_j\varphi(W_B)\right]
 \longrightarrow
 \mathbb E[\mathbf1_{A_0}W_B\varphi(W_B)].
\]
The limit uses the $L^2$ convergence defining $W_B$; adding index $n$
does not change that frequency.  Taking
$\varphi(x)=(1+zx)^{-n}$ and using $W_B\le W$ on $A$ now gives
\begin{align}
 \mathbb E[\mathbf1_A(1+zW)^{-n}]
 &\leq\mathbb E[\mathbf1_A(1+zW_B)^{-n}]\notag\\
 &=\mathbb E[\mathbf1_{A_0}W_B(1+zW_B)^{-n}]
 \leq z^{-1}.
 \label{arr:duplication-contradiction}
\end{align}
The last bound uses $x/(1+zx)^n\leq1/z$.  Equations
\eqref{arr:duplication-transform} and
\eqref{arr:duplication-contradiction} contradict $b<1$ as $z\to\infty$.
Consequently the desired extension event has positive probability.
Shrinking the neighborhoods and decreasing $b_*$ to its target value,
compactness of the bounded finite-block support proves
\eqref{arr:duplication}.  Permutation symmetry permits duplication of
any chosen index.

If ultrametricity failed, a supported triple could be chosen with
$a_{12}<a_{13}\leq a_{23}$.  Repeated duplication makes three groups of
$m$ indices, preserving the three cross-group entries $a_{12},a_{13},a_{23}$.
Within each group all off-diagonal entries stay at most its largest
cross-group entry: this follows inductively from
\eqref{arr:duplication}.  Represent the resulting finite Gram matrix by
vectors and let $u_1,u_2,u_3$ be their group averages.  Then
\[
 \|u_2\|^2,\|u_3\|^2
 \leq \frac Dm+\left(1-\frac1m\right)a_{23},
 \qquad u_2\cdot u_3=a_{23},\qquad \|u_1\|\leq\sqrt D.
\]
Hence $\|u_2-u_3\|^2\leq2(D-a_{23})/m\to0$, whereas
$u_1\cdot(u_3-u_2)=a_{13}-a_{12}>0$.  Cauchy--Schwarz gives a
contradiction.  Applying exchangeability to all triples completes the
proof.
\end{proof}

\subsection{Finite cascades and their sampling law}

The hierarchical weights defined next are finite Ruelle probability
cascades~\cite{ruelle1987}. Their complete weights-and-ancestry sampling
law was identified by Bolthausen and Sznitman~\cite{bolthausensznitman1998}.
We derive the needed sampling and transformation laws from the Poisson
construction, including repeated visits to the same leaf.

Fix $0=\zeta_0<\zeta_1<\cdots<\zeta_k=1$ and put
$w_i=\zeta_{i+1}-\zeta_i$.  On a rooted tree of depth $k-1$, the outgoing
edge weights at each vertex of depth $i-1$ are the points of an independent
Poisson process on $(0,\infty)$ with intensity
\[
 \zeta_i x^{-1-\zeta_i}\,dx,\qquad 1\leq i\leq k-1.
\]
Children may be indexed in decreasing order of these points.  If $u_\gamma$
is the product of the edge weights on the path to leaf $\gamma$, define
\begin{equation}
 T=\sum_\gamma u_\gamma,
 \qquad v_\gamma=u_\gamma/T.
 \label{arr:cascade-weights}
\end{equation}
For depth zero there is one leaf and $T=v_\gamma=1$.  Write
$\gamma\wedge\eta$ for the depth of the deepest common vertex of two
leaves; it ranges from $0$ to $k-1$, including when the leaves are equal.
The \emph{shape} of sampled labels records their nested ancestral
partitions, without the numerical child indices.

\begin{lemma}[Cascade law and finite-level uniqueness]
\label{arr:cascade}
The total $T$ in \eqref{arr:cascade-weights} is almost surely finite and
positive, with $\mathbb E[(\log T)^2]<\infty$.  For replicas sampled from
$(v_\gamma)$, the common-level array satisfies the Ghirlanda--Guerra
property and
\begin{equation}
 \mathbb P(\gamma^1\wedge\gamma^2=i)=w_i,
 \qquad 0\leq i\leq k-1.
 \label{arr:level-marginal}
\end{equation}
Conversely, a weakly exchangeable ultrametric scalar array with the
Ghirlanda--Guerra property and finitely many off-diagonal values has the
same off-diagonal law as this cascade, with the levels replaced by its
ordered values and with $w_i$ their marginal probabilities.  Values of
zero probability are omitted.
\end{lemma}

\begin{proof}
We give the Poisson and sampling calculations, since the distribution of a
descendant tree under sampling requires a power bias.  If a Poisson process
has intensity $C x^{-1-b}\,dx$, $0<b<1$, and its points are multiplied by
independent positive marks $X$ with $0<\mathbb E X^b<\infty$, the image
process has intensity $C\mathbb E X^b\,x^{-1-b}\,dx$.  More precisely,
after this change the attached mark law is biased by $X^b$, independently
of the image point.  This follows by substituting $y=xX$ in the marked
intensity.  Independent marking and mapping follow directly from the
Poisson count law on disjoint sets, first for simple marks and then by
approximation.

The sum $S$ of points of intensity $C x^{-1-b}\,dx$ has Laplace transform
\begin{equation}
 \mathbb E e^{-tS}=e^{-ct^b},
 \qquad c=C\Gamma(1-b)/b>0.
 \label{arr:stable-laplace}
\end{equation}
Indeed the Poisson exponential formula gives the exponent
$-C\int_0^\infty(1-e^{-tx})x^{-1-b}\,dx=-ct^b$.
The integral is finite; the formula proves both finiteness and positivity
of the sum.  For $0<r<b$, integrating
$1-\mathbb E e^{-tS}$ against $t^{-1-r}\,dt$ gives
$\mathbb E S^r<\infty$.  For every $r>0$, the identity
\[
 \mathbb E S^{-r}
 =\frac1{\Gamma(r)}\int_0^\infty t^{r-1}e^{-ct^b}\,dt
\]
gives a finite negative moment.  These moments imply a finite second
moment for $\log S$.  Summing the tree from its last level upward and
using the marking calculation proves the assertions about $T$: the
descendant exponent is strictly larger than the exponent of the edge
above it, so every required positive fractional moment exists.

At a parent, let $b$ be the exponent on its outgoing edges.  Multiply
each edge point by its descendant total.  The resulting child masses are
Poisson points of intensity $C x^{-1-b}\,dx$, and their attached normalized
descendant trees are independent of the points and have law biased by the
descendant total to power $b$.  If the parent tree itself is biased by its
total to a power $\lambda<b$, only the point process is further biased by
the $\lambda$th power of its sum; the conditional descendant description
does not change.  At the root $\lambda=0$, and thereafter $\lambda$ is
the incoming exponent.

Here is the resulting partition law.  Denote unbiased child masses by
$(X_j)$ and their sum by $S$.  Under the bias $S^\lambda$, the probability
of a specified partition of $m\geq1$ visits into $d$ nonempty child classes
of sizes $m_1,\ldots,m_d$ is
\begin{align}
 p_\lambda(m_1,\ldots,m_d)
 &=\frac{\mathbb E\left[S^{-m+\lambda}
          \sum_{j_1,\ldots,j_d\ \mathrm{distinct}}
                     \prod_{r=1}^dX_{j_r}^{m_r}\right]}
         {\mathbb E S^\lambda}\notag\\
 &=\frac{\prod_{r=1}^d C\Gamma(m_r-b)}
        {\Gamma(m-\lambda)\mathbb E S^\lambda}
   \int_0^\infty t^{db-\lambda-1}e^{-ct^b}\,dt.
 \label{arr:partition-integral}
\end{align}
To derive this, insert the Laplace integral for $S^{-m+\lambda}$ and use
the Poisson exponential formula with $d$ distinct point insertions.  Each
insertion contributes
$C\int_0^\infty x^{m_r-b-1}e^{-tx}\,dx
=C\Gamma(m_r-b)t^{b-m_r}$.
For completeness, the insertion rule follows by first restricting to a
finite-intensity process, summing its Poisson count series, and then
letting the truncation grow; Tonelli's theorem applies to the nonnegative
integrands.

The ratio of \eqref{arr:partition-integral} after and before a visit to
an existing class $r$ is
\begin{equation}
 \frac{m_r-b}{m-\lambda}.
 \label{arr:join-child}
\end{equation}
The complementary probability of creating a new child is
$(db-\lambda)/(m-\lambda)$.  Observed deeper shapes do not alter
\eqref{arr:join-child}, because the attached descendant trees are
independent of the child masses.  Equivalently the probabilities of
complete shapes factor into the successive partition probabilities.
Along an existing class at level $i$, these ratios telescope: after $n$
replicas, the probability that the next replica enters that class is
\begin{equation}
 \frac{m_*-\zeta_i}{n},
 \label{arr:join-level}
\end{equation}
where $m_*$ is its size.  At level zero the class is the whole tree and
the probability is one.  With one previous replica,
\eqref{arr:join-level} gives
$\mathbb P(\gamma^1\wedge\gamma^2\geq i)=1-\zeta_i$, proving
\eqref{arr:level-marginal}.  With $n$ previous replicas, applying
\eqref{arr:join-level} to the class of replica $1$ gives exactly
\eqref{arr:conditional-gg}, first for upper-threshold tests and hence for
every test on the finite set of levels.

Conversely, the upper-threshold relations of a finite ultrametric array
are nested equivalence relations.  Conditional on its first $n$ indices,
the Ghirlanda--Guerra property gives
\eqref{arr:join-level} for each existing class, using a representative of
that class.  These probabilities specify the distribution of the deepest
existing class joined by the new index, and therefore of all its new
pair entries.  They coincide with the cascade probabilities.  Induction
on $n$, starting from the prescribed two-replica law, proves uniqueness.
\end{proof}

\subsection{Quantization and continuity of array laws}

\begin{lemma}[Ordered quantization with the original diagonal]
\label{arr:quantization}
Let $B$ be nonnegative and ultrametric off the diagonal, with constant
diagonal $D$ and entries at most $D$.  Applying any nonnegative
nondecreasing finite-valued map to its off-diagonal entries, with image
contained in $[0,D]$, and retaining diagonal $D$, produces a Gram array.
If $B$ has the Ghirlanda--Guerra property, its quantized off-diagonal law
is the finite cascade law in Lemma~\ref{arr:cascade}.

More generally, suppose finitely many desired off-diagonal covariance
entries are bounded nonnegative nondecreasing functions $a_j(r)$ of the
scalar entries $r=B_{\ell m}$, $\ell\ne m$, with $a_j(r)\leq D_j$.
One can choose finite ordered
bins of $r$ that approximate $r$ and every $a_j(r)$ uniformly on their
domain, with errors tending to zero, while keeping each covariance
diagonal equal to $D_j$.  All the resulting covariance arrays are Gram.
The binning can be chosen as a function of $r$ alone, so it never splits
an atom of the scalar overlap law.
\end{lemma}

\begin{proof}
Let the distinct quantized values be $c_0<\cdots<c_s$.  For a finite
replica restriction, let $A_j$ be the matrix indicating that two indices
belong to the same threshold class at level $c_j$, with diagonal entries
one.  Every $A_j$ is positive semidefinite: its quadratic form is the sum
over classes of the squared sum of the coordinates in that class.  The
quantized matrix is
\begin{equation}
 c_0\mathbf1\mathbf1^{\mathsf T}
 +\sum_{j=1}^s(c_j-c_{j-1})A_j+(D-c_s)I.
 \label{arr:gram-decomposition}
\end{equation}
All coefficients are nonnegative.  The Ghirlanda--Guerra property passes
through the quantization map by the tower rule, so
Lemma~\ref{arr:cascade} identifies its law.

For the joint assertion, bin $r$ and each $a_j(r)$ on meshes of size
$\varepsilon$ and take their common refinement.  Since every coordinate
is nondecreasing, the nonempty joint bins are totally ordered.  Choosing
one actual vector $(r,a_1(r),\ldots,a_d(r))$ in each bin gives ordered
representatives and error at most $\varepsilon$ in each coordinate.
Use \eqref{arr:gram-decomposition} separately for each coordinate, with
its own diagonal $D_j$.  The bins depend only on $r$, including at atoms.
If the functions are specified only on a full-measure ordered domain,
the construction on that domain suffices for every sampled array entry.
\end{proof}

\begin{lemma}[Uniqueness and bounded-Lipschitz continuity]
\label{arr:law-continuity}
For fixed diagonal $D$, the law of a nonnegative ultrametric
Ghirlanda--Guerra array with values in $[0,D]$ is determined by its
two-replica marginal.  Such a law exists for every probability marginal
on $[0,D]$.  If the marginal quantiles $q_m,q$ satisfy
$\|q_m-q\|_{L^1(0,1)}\to0$, then the laws of every fixed finite replica
block converge in bounded-Lipschitz distance.
\end{lemma}

\begin{proof}
For uniqueness, round two candidate arrays on increasingly fine finite
meshes.  Their rounded marginals coincide, so
Lemma~\ref{arr:cascade} gives identical rounded array laws.  Uniformly
vanishing rounding error then identifies their original finite-block
laws.  For existence, approximate the desired quantile by ordered step
functions, use finite cascades, and extract a limit of the bounded
arrays.  Their Gram property follows from
\eqref{arr:gram-decomposition}; ultrametricity and the
Ghirlanda--Guerra property pass to the limit.  The limit has the desired
marginal.

If $q_m\to q$ in $L^1$, their marginals converge weakly (in fact, coupling
them by the same uniform variable bounds their Wasserstein distance by
$\|q_m-q\|_1$).  Any subsequential limit of the array laws has the
Ghirlanda--Guerra property, is nonnegative and ultrametric, and has this
limiting marginal.  Uniqueness identifies it, proving convergence of the
whole sequence.  Weak convergence on the compact space of finite blocks
is equivalent to convergence in bounded-Lipschitz distance.  This is a
continuity assertion; no rate independent of the finite-block size is
needed.
\end{proof}

\subsection{Poisson displacement and the tilted mark law}

The overlap sampling law does not by itself determine observables of the
fresh Gaussian fields after reweighting.  For example, for smooth $f$,
restoring a pattern also changes the sampled values of $f'$.  We therefore need
the joint transformation of cascade weights and marks, developed by
Panchenko and Talagrand~\cite[Lemmas~2.1 and~3.1]{panchenkotalagrand2007property};
see also \cite[Section~3]{panchenkotalagrand2007interpolation} for its role in
Guerra's interpolation.  The following form keeps track of the root law
and of independent additive mark components.

Draw root data independently of all Poisson edge processes.  Conditional
on an ancestor's marked path, give its children independent marks with a
common probability kernel, independently of the edge points.  These
kernels may depend on the marked path and its depth, but not on numerical
child indices or Poisson weights.
A shared Gaussian field with covariance levels
$0\leq b_0\leq\cdots\leq b_{k-1}$ is obtained by a common root increment
of variance $b_0$ and independent vertex increments of variances
$b_i-b_{i-1}$.

\begin{lemma}[Marked cascade transformation]
\label{arr:tilt}
Let $X_{k-1}$ be a terminal log factor depending on the marked path.
For $i\ge1$, let $\mathbb E_i$ integrate the original level-$i$
child-mark kernel, conditional on the marked ancestor path through level
$i-1$, without averaging the Poisson weights or root data. Define
\begin{equation}
 X_{i-1}=\frac1{\zeta_i}\log\mathbb E_i e^{\zeta_i X_i},
 \qquad i=k-1,\ldots,1.
 \label{arr:mark-recursion}
\end{equation}
Assume these transforms are finite; the bounded log factors and the
Gaussian log factors with at most linear growth used below satisfy this
condition.  For
$Z=\sum_\gamma v_\gamma e^{X_{k-1}(\gamma)}$,
\begin{equation}
 \mathbb E[\log Z\mid\text{root data}]=X_0,
 \qquad
 \mathbb E[(\log Z-X_0)^2\mid\text{root data}]
 \leq4\mathbb E[(\log T)^2].
 \label{arr:log-cascade}
\end{equation}
Here $T$ is the unmarked total in \eqref{arr:cascade-weights}.

After reweighting by $e^{X_{k-1}}$ and relabeling children, the joint law
of normalized weights and marks can be generated by first drawing an
ordinary unmarked cascade, independently of the root data, and then
drawing marks on its vertices with kernels of densities
\begin{equation}
 e^{\zeta_i(X_i-X_{i-1})},\qquad 1\leq i\leq k-1,
 \label{arr:tilted-kernel}
\end{equation}
relative to the original kernels, with independent draws on distinct
children conditional on ancestor marks.  Root data retain their original
law.  In particular, sampling reweighted leaves preserves their shape law;
conditioning on that shape leaves exactly these mark kernels along the
sampled branches.  Independent mark components remain independent under
these kernels when their original kernels factor and the terminal log
factor is a sum of the separate component log factors.
\end{lemma}

\begin{proof}
At every edge from level $i-1$ to $i$, multiply its Poisson point by
$e^{X_i-X_{i-1}}$ and reorder children.  By
\eqref{arr:mark-recursion}, the conditional $\zeta_i$th moment of this
multiplier is one.  The marked Poisson mapping calculation in the proof
of Lemma~\ref{arr:cascade} therefore leaves the point intensity unchanged
and changes the attached mark kernel exactly by
\eqref{arr:tilted-kernel}.

The independence assertion requires doing this from the last level
upward.  Inductively, conditional on the path to a child, its transformed
unmarked subtree has the ordinary unmarked law, independently of that
path; its marks follow the transformed descendant kernels.  Attach the
entire transformed subtree as part of the child's mark in the mapping
at its parent.  The multiplier depends only on the path through that
child, so the mapping tilts the child kernel while keeping this
conditional subtree law.  This proves at the root that the transformed
unmarked tree has its original law independently of root data, with
marks generated by the stated conditional kernels.

Products of the edge multipliers telescope.  If $T^*$ is the transformed
raw total, then
\begin{equation}
 Z=e^{X_0}\frac{T^*}{T}.
 \label{arr:raw-total-ratio}
\end{equation}
Both $T$ and $T^*$ have the same conditional law given root data, with
finite squared logarithms by Lemma~\ref{arr:cascade}.  Taking the
conditional expectation of the logarithm proves the first part of
\eqref{arr:log-cascade}; the elementary square bound proves the second.
No independence between $T$ and $T^*$ is asserted or required.

The reweighted normalized masses are precisely the transformed masses,
up to a relabeling that preserves ancestry.  Sampling them consequently
gives the ordinary cascade shape law.  Conditional on the transformed
unmarked tree, its mark kernels have the description just proved; hence
sampling or conditioning on a shape does not further tilt those kernels.
This independence concerns the transformed tree, not the original
weights.  Finally, factoring the expectations in
\eqref{arr:mark-recursion} for an additive terminal log factor shows that
the transforms add and the tilted component kernels factor.
\end{proof}

\paragraph{Terminal variance and repeated visits.}
Suppose a Gaussian mark has true diagonal variance $D$ and shared
off-diagonal levels $b_i$, with $b_{k-1}\leq D$.  The cascade represents
the shared field $y_\gamma$.  On each replica visit introduce an
independent residual Gaussian of variance $D-b_{k-1}$.  Thus, for a
positive leaf factor $A$, the factor to use in the cascade recursion is
\begin{equation}
 \overline A(y)=\mathbb E_\eta
   A\bigl(y+\sqrt{D-b_{k-1}}\,\eta\bigr).
 \label{arr:residual-average}
\end{equation}
After reweighting, a visited residual has density
$A(y+\sqrt{D-b_{k-1}}\eta)/\overline A(y)$ relative to its standard
Gaussian law, independently on different visits conditional on their
shared marks.  This is an ordinary coefficient-one tilt.  In particular,
two replicas visiting the same leaf still use two independent residuals.
Giving a leaf one permanent residual would incorrectly raise their
off-diagonal covariance from $b_{k-1}$ to $D$.  These statements also hold
for independent vector residuals and for additional finite-spin variables
integrated in the leaf factor.  The residual prescription realizes
exactly the final diagonal term in \eqref{arr:gram-decomposition}.

\begin{figure}[tb]
 \centering
 \begin{tikzpicture}[>=Stealth, every node/.style={font=\small},
                     shared/.style={draw,rounded corners,inner sep=5pt},
                     visit/.style={draw,dashed,rounded corners,inner sep=5pt}]
  \node[shared] (root) at (0,0) {common root $Z_0$};
  \node[shared] (left) at (-2,-1.25) {leaf $\gamma$: $y_\gamma$};
  \node[shared] (right) at (2,-1.25) {leaf $\eta$: $y_\eta$};
  \node[visit] (vone) at (-3.3,-2.6) {visit $1$: $y_\gamma+\epsilon_1$};
  \node[visit] (vtwo) at (-.1,-2.6) {visit $2$: $y_\gamma+\epsilon_2$};
  \node[visit] (vthree) at (3.3,-2.6) {visit $3$: $y_\eta+\epsilon_3$};
  \draw[->] (root)--(left);
  \draw[->] (root)--(right);
  \draw[->,dashed] (left)--(vone);
  \draw[->,dashed] (left)--(vtwo);
  \draw[->,dashed] (right)--(vthree);
 \end{tikzpicture}
 \caption{A two-leaf schematic of the Gaussian construction before
 reweighting. Shared increments
 belong to tree vertices, while residual noises belong to replica visits.
 The centered Gaussian residuals $\epsilon_1,\epsilon_2,\epsilon_3$ are
 independent and have variance $D-b_{k-1}$, even though visits $1$ and $2$
 have the same leaf label.}
 \label{arr:residual-figure}
\end{figure}

\subsection{Fresh Gaussian fields and passage through array limits}

\begin{lemma}[Bounded reweighting of convergent arrays]
\label{arr:limit}
Let a sequence of random base probabilities have jointly convergent
arrays of bounded observables and of the conditional covariances of
finitely many fresh centered Gaussian fields.  The fields are sampled
independently of the base probability conditional on its data, with any
specified independence between field components.  Then the following
averages converge, as determined by those limiting arrays:
\begin{enumerate}
 \item bounded continuous functions of the listed observable and field
 evaluations on finitely many base replicas;
 \item the same tests after reweighting by a fixed continuous function
 of the listed evaluations at one replica, bounded between two fixed
 positive constants;
 \item the expected logarithm of the normalizing integral of that factor.
\end{enumerate}
Independent finite-spin averages may be included.  The assertion remains
valid at singular limiting covariance matrices.  If the limiting
off-diagonal covariance entries are ordered functions of a scalar
ultrametric Ghirlanda--Guerra array, its evaluations can be obtained by
the quantized cascades of Lemma~\ref{arr:quantization}, retaining the true
diagonal variances and using \eqref{arr:residual-average}, provided the
base observables and tests used in the evaluation are also approximated
by these bins in probability on every finite replica restriction.  This
condition includes constants and continuous tests of the scalar overlap
and the quantized covariance coordinates.
\end{lemma}

\begin{proof}
On any fixed number of base replicas, the conditional joint Gaussian
law is a continuous function of its covariance matrix, also on the
boundary of the positive semidefinite cone.  For example, couple it as
$C^{1/2}g$; the nonnegative matrix square root is continuous and the
Gaussian vector $g$ has fixed dimension.  Bounded continuous tests
therefore pass to the joint array limit.

Write the positive reweighting factor as $A$, and let $Z=P(A)$ for the
base probability $P$.  If $0<c\leq A\leq C$, then $c\leq Z\leq C$.
An $r$-replica reweighted test has expectation
\[
 \mathbb E\left[
   Z^{-r}P^{\otimes r}\left(F\prod_{\ell=1}^r A_\ell\right)
 \right].
\]
Uniformly approximate $z^{-r}$ on $[c,C]$ by polynomials.  Each monomial
$Z^j$ is represented by $j$ additional base replicas, evaluated in the
same fresh fields.  The resulting terms are finite-replica tests covered
by the first part, and the uniform approximation controls the error.
Approximating $\log z$ on the same interval proves the logarithmic
assertion.  Finite-spin variables are simply additional independent base
variables in this argument.

Finally quantize all the ordered covariance coordinates jointly.  Their
finite-replica matrices converge by
Lemma~\ref{arr:quantization}, and the other observables and tests converge
by the additional approximation hypothesis.  Gaussian continuity and
the same polynomial argument therefore apply.  The shared Gaussian construction with
independent residuals has exactly those matrices, including on repeated
leaf visits.  Lemma~\ref{arr:cascade} identifies the quantized shape law,
and Lemma~\ref{arr:tilt} evaluates its bounded reweightings.  Unbounded
factors require an additional uniform-integrability argument; the cavity
calculation below supplies the particular truncation estimate it needs.
\end{proof}

\begin{lemma}[Evaluation of one fresh pattern]
\label{arr:fresh-pattern}
Suppose that the limiting spin-overlap array of a sequence of base Gibbs
probabilities has diagonal one, is Gram, and satisfies the
Ghirlanda--Guerra property.  Let $p$ be the quantile of its off-diagonal
two-replica law.  For a fresh independent Gaussian pattern field $g(x)$
with conditional covariance $R(x,x')$, and bounded continuous $f$,
\begin{equation}
 \lim\mathbb E\log\langle e^{f(g(x))}\rangle=V_f(p),
 \label{arr:fresh-pattern-value}
\end{equation}
whenever $V_f$ is understood through its continuous quantile extension.
In particular this applies to the smooth potentials used in the proof
before the final approximation step.
\end{lemma}

\begin{proof}
Theorem~\ref{arr:ultrametric} gives positivity and ultrametricity.
Lemma~\ref{arr:limit} passes the bounded positive factor through the array
limit.  Round the off-diagonal overlap into levels
$p_0\leq\cdots\leq p_{k-1}$ with probabilities $w_i$, leaving its
diagonal one.  On the corresponding cascade, first integrate the
independent variance $1-p_{k-1}$ at each visit.  The terminal shared-field
log factor is consequently
\[
 U_{k-1}(x)=\log\mathbb E
      e^{f(x+\sqrt{1-p_{k-1}}G)}.
\]
The remaining branch transforms have coefficients
$\zeta_{k-1},\ldots,\zeta_1$, by Lemma~\ref{arr:tilt}, and the common
root Gaussian is averaged with coefficient zero.  These are precisely
the defining transforms of $V_f$ on the rounded step quantile: its final
coefficient is one, and its initial coefficient is zero.  Sending the
rounding error to zero proves \eqref{arr:fresh-pattern-value}, by
Lemma~\ref{arr:limit} and continuity of the recursion in its quantile.
For bounded continuous $f$ this also supplies a direct description of
that extension: Lemma~\ref{arr:law-continuity} and the bounded reweighting
argument show that the cascade values have a unique continuous limit.
\end{proof}

\subsection{Gaussian Gibbs calculus and variance bounds}

\begin{lemma}[Gaussian differentiation]
\label{arr:calculus}
Let $P$ be a probability reference measure, possibly random but
independent of the Gaussian fields under consideration.  For independent
centered Gaussian energies with covariances $C_0,C_1$, interpolate using
$H_t=\sqrt{1-t}\,H_0+\sqrt t\,H_1$, and put
$\Delta=C_1-C_0$.  Then
\begin{equation}
 \frac{d}{dt}\mathbb E\log P(e^{H_t})
 =\frac12\mathbb E\langle\Delta_{11}-\Delta_{12}\rangle_t.
 \label{arr:gaussian-log-derivative}
\end{equation}
For a bounded $r$-replica test $F$ independent of these fields,
\begin{align}
 \frac{d}{dt}\mathbb E\langle F\rangle_t
 =\frac12\mathbb E\Bigg\langle F\Bigg(
 &\sum_{a,b=1}^r\Delta_{ab}
 -2r\sum_{a=1}^r\Delta_{a,r+1}\notag\\
 &+r(r+1)\Delta_{r+1,r+2}-r\Delta_{r+1,r+1}
 \Bigg)\Bigg\rangle_t,
 \label{arr:replica-derivative}
\end{align}
so its magnitude is at most
$(2r^2+r)\|F\|_\infty\|\Delta\|_\infty$.
If an independent centered Gaussian field $Y$ of covariance $C$ enters
the energy with coefficient $s$, Gaussian integration by parts also gives
\begin{equation}
 \mathbb E\langle Y(1)F\rangle
 =s\,\mathbb E\left\langle
  F\left(\sum_{a=1}^r C_{1a}-rC_{1,r+1}\right)
 \right\rangle.
 \label{arr:gaussian-insertion}
\end{equation}
These formulas hold for finite spin spaces with countably many cascade
labels whenever the covariances are bounded at the fixed system size.
They can be integrated up to the endpoints of covariance segments.
\end{lemma}

\begin{proof}
For finitely many states, differentiate the coupling $H_t$ and apply
ordinary Gaussian integration by parts.  This contracts one half of the
energy Hessian with $\Delta$.  A derivative of a normalized $r$-replica
integral inserts the sum over its $r$ sampled states minus $r$ times the
insertion at one extra replica.  A second derivative and contraction
give \eqref{arr:replica-derivative}; the sum of its absolute coefficients,
including the factor $1/2$, is $2r^2+r$.  The first derivative of the log
partition gives \eqref{arr:gaussian-log-derivative}, and a single
integration by parts gives \eqref{arr:gaussian-insertion}.

To pass to countable labels, truncate to increasing finite label sets
chosen independently of the added fields and renormalize their reference
weights.  Uniform integrability follows from elementary Gaussian moment
bounds.  For example, if a Gaussian energy has pointwise variance at most
$K$, then, conditionally on its independent probability reference,
\begin{equation}
 \mathbb E Z^p\leq e^{p^2K/2},\qquad
 \mathbb E Z^{-p}\leq e^{p^2K/2},\qquad
 Z=P(e^H),\quad p\geq1.
 \label{arr:partition-moments}
\end{equation}
The first inequality is Jensen's inequality for $x^p$; the second uses
$[P(e^H)]^{-p}\leq P(e^{-pH})$.  Cauchy--Schwarz and the same
single-state Gaussian moments control polynomial Gaussian insertions,
and exponentials of sums of absolute values of finitely many such
fields.  Thus Gibbs derivatives, their Gaussian insertions, and the
logarithms have uniformly integrable bounds on compact coefficient sets.
The truncated integrals converge almost surely and in the necessary
moments.  Bounded pattern log factors can be absorbed into the reference;
if estimates retain them explicitly, their exponential costs are
integrable for either a fixed count or an independent Poisson count.

Integrate the finite-state covariance derivative first and then pass to
the limit using its uniform bound.  Gaussian coupling and the same
moment bounds give continuity at segment endpoints.  Ordinary energy
parameter derivatives and \eqref{arr:gaussian-insertion} pass to the
limit by the polynomial-insertion bounds.  This proves the stated
extensions.
\end{proof}

\begin{lemma}[Elementary variance estimates]
\label{arr:variance}
A function of finite standard Gaussian data with Euclidean Lipschitz
constant $A$ has variance at most $A^2$.  A function of independent
inputs whose value changes by at most $c_i$ when input $i$ is replaced
has variance at most $\sum_i c_i^2$.  If a function $a$ on the integers
has Lipschitz constant $c$ and $K$ is Poisson with mean $\lambda$, then
$\operatorname{Var}(a(K))\leq c^2\lambda$.
\end{lemma}

\begin{proof}
For a smooth function of Gaussian data, interpolate two standard
Gaussian copies from independent to identical, with cross covariance
$tI$.  Gaussian integration by parts differentiates their product
expectation into the expectation of their gradient inner product.
Integrating from zero to one and using the gradient bound proves the
first claim.  Truncation and smoothing extend it to Lipschitz functions.
For the second, reveal the independent inputs successively and sum the
variances of the martingale differences, each bounded by $c_i^2$.
For the third, let $K'$ be an independent copy and use
\[
 \operatorname{Var}(a(K))
 =\tfrac12\mathbb E[(a(K)-a(K'))^2]
 \leq\tfrac{c^2}{2}\mathbb E[(K-K')^2]=c^2\lambda.
\]
\end{proof}
